Project Kennel is a userspace reference monitor for the interactive user
level, a way to run code under your own account without granting it your
own authority. This is the second of its two volumes; the first set out
the design without reference to any operating system, and this one
realises that design on Linux. The altitude is inverted on purpose.
Where the first volume held above the mechanism, refusing to name a
syscall or a kernel feature so that its claims would read the same on
any port, this one names them throughout: the kernel facilities the
confinement is woven from, the processes that hold and drop privilege,
the bus the workload talks across, the crates the trusted base is built
in. Every claim here answers a question the first volume left open on
purpose, the question of how, on this kernel, a contract stated there is
kept.

The volume is in four parts. The first is the machinery: the kernel
facilities and how they are woven into a confinement, the processes that
build and supervise a kennel and the single one that holds privilege,
the binder bus every kennel talks across, the crate topology of the
trusted base, and the command surface that drives it. This part has no
analogue in the first volume, because it is the substrate the first
volume deliberately did not name. The second takes the resource classes
the first volume's second part defined and gives each its Linux
mechanism, from the filesystem through execution and the network to the
local, cryptographic, and image surfaces and the smaller remaining ones.
The third is the interplay: how a kennel spawns another, how kennels
offer and consume one another's capabilities, and how a service too rich
to broker is interposed as a confined instance of itself. The fourth is
the language and the tools: the grammar a confinement is written in, the
compiler that turns it into the artefact the runtime enforces, the keys
that sign it, and the account the monitor keeps. The first volume did
not separate the last two: it carried the interplay and the declared
grant together, and this volume takes them apart because the mechanism
does.

Three citation forms point out of the prose, into two registers and a
catalogue. A name such as \texttt{construction-by-absence} points into
the first volume's principle register (\texttt{PRINCIPLES.md}): it names
the design posture a mechanism keeps, the contract this volume is here
to honour. A name such as \texttt{quarantine-the-unsafe} points into
this volume's own register (\texttt{CONSTRUCTION.md}): it names a
construction principle the code itself holds to, a property checkable by
opening the repository rather than by reading the design. An identifier
such as \texttt{T5.1} points into the threat catalogue
(\texttt{THREATS.md}), shared with the first volume, and names a
concrete threat the mechanism in hand closes. Literature is cited by
name in brackets and resolved at the end, as in the first volume.

This volume assumes the first. A reader can follow the mechanism here
without the design behind it, but the reason a mechanism is shaped the
way it is lives in the first volume, and the principle citations are the
way back to it. The machinery of the first part is assumed by every
chapter after it, so it is the one part meant to be read before the
others; a resource chapter then stands largely on its own, with its
citations leading back to the design it realises and the threats it
answers. The kernel floor is the project floor throughout, kernel 6.10
for Landlock ABI 5, and a mechanism that needs more says so where it is
named.
